

% Placeholder marker -- fill these in and then delete the \ph wrapper.
\newcommand{\ph}[1]{\textcolor{red}{[#1]}}

\chapter{Methodology}

\section{Overview of the approach}\label{sec:overview}

The goal of this work is a solid, thickness-accurate 3D model of the valve
leaflets that can serve as the basis for an ultrasound phantom. The starting
point is two things: a volumetric ultrasound scan of the valve, stored as a VDB
volume, and a surface mesh of the leaflets. The problem is that these two do not
directly provide what is needed. The scan contains the full valve, but only as a
grid of grayscale values with no explicit geometry. The mesh has clean geometry,
but it is only a single surface with no thickness, so on its own it cannot be
turned into a physical object with a wall.

The core task of the whole pipeline is therefore to bridge these two: to bring
the mesh and the volume into a common coordinate system, to recover the leaflet
thickness from the volume, to transfer that thickness to the mesh, and to turn
the single surface into a closed solid.

I did not carry out all of this in a single program, because no single tool does
every part well. Instead the work is split across three: MeVisLab is used to load
and prepare the volume, 3D~Slicer is used to align the mesh with the volume and
to take the thickness measurements, and Blender is used to interpolate the
measurements over the mesh and build the solid. Some smaller steps, mainly the
interpolation, were done with short Python scripts using the numpy librariy. The following sections describe each part, and, more importantly, why
each choice was made and not one of the other possible ones.

\section{Input data}\label{sec:input}

The first input is the volumetric scan. It is a 3D grayscale field where bright
voxels roughly correspond to tissue and dark voxels to the surrounding space. It
contains the leaflets with their real thickness, but this thickness is only
present as a bright band in the image and not as a number or a shape that can be
measured automatically without further steps. The scan is also not clean
everywhere: the ultrasound probe used to record it left a very noisy region in
the centre of the valve, around the opening, which becomes important later.

The second input is the surface mesh of the leaflets. This mesh is a single
surface. It does not enclose a volume, and it has no wall thickness at all. In
practice there are two separate leaflet meshes, referred to here as \ph{Mesh1}
and \ph{Mesh2}.

Before deciding how to add thickness, I first had to establish what the surface
actually represents, because this changes the whole calculation. A surface mesh
of a thin structure can either be a \emph{midsurface}, sitting in the middle
between the two walls, or a \emph{boundary surface}, sitting on one side. If it
is a midsurface, the real thickness is split evenly on both sides of it. If it is
a boundary surface, the whole thickness sits on one side only. Getting this wrong
would make every thickness value off by a factor of two.

To check this, the mesh was placed next to the volume and a cross-section through
a leaflet was examined. If the mesh were a midsurface, it would run through the
middle of the bright leaflet band. In this case the surface sits on one side of
the band, which means it is the underside (a boundary surface), not the middle.
This has a direct consequence for the later steps: the full measured thickness
has to be added on one side, and there is no division by two. I come back to this
in Section~\ref{sec:solid}.

\section{Choosing how to recover thickness}\label{sec:recover}

Thickness is not stored anywhere in the data. It has to be recovered from the
grayscale volume. There were two reasonable ways to do this, and it is worth
explaining both, because the reason for choosing the second one is part of the
method.

The first idea was to compute the thickness fully automatically with a distance
transform. The steps would be: turn the volume into a binary mask, where material
is~1 and everything else is~0, and then run a distance transform on that mask. A
distance transform gives, for every inside voxel, the distance to the nearest
background voxel, in other words the distance to the closest wall \cite{edMorphologyDistance}. In the middle
of a leaflet, that distance is half the local thickness, so the thickness could
in principle be read off as two times the distance value. This was set up in
MeVisLab using an \texttt{OpenVDBLoad} module to load the volume, a
\texttt{Threshold} module to make the binary mask (values above~0.4 set to~1,
everything else to~0), and a \texttt{EuclideanDistanceTransform} module in 3D
mode to produce the distance field.

The reason this idea is attractive is that it is automatic and dense: it would
give a thickness value everywhere on the leaflet at once, without any manual
work. In the end, though, I decided not to use it, for a few reasons that add up.

First, there is a practical problem with reading the result. To get thickness at
the mesh, the distance field has to be sampled exactly at the positions of the
mesh points. MeVisLab does not have a simple, ready-made module for ``sample this
volume at the vertices of that mesh'', and the scripting module has no image
input or output that would let this be done directly inside the network. It would
mean exporting the distance field and doing the sampling externally in Python,
which is possible but adds another step that would itself have to be checked.

Second, the simple ``two times the distance'' rule only holds if the mesh is a
midsurface. As shown in Section~\ref{sec:input}, the mesh here is the underside.
For a boundary surface the correct value would be the distance from the underside
across to the top wall, which is not what a symmetric distance transform gives
without extra work. So the clean version of the method does not actually match
this data.

Third, and most importantly, the whole result depends on one global threshold
that separates material from background, and the volume is not clean enough for
that. The Ultrasound volume has speckled and blurred edges, and this scan has an
especially noisy region in the centre of the valve. A single fixed threshold
applied to the whole volume would turn that noise into false geometry exactly
where the data is worst, and there would be no easy way to notice that it had
happened. A dense automatic field of numbers is only useful if it can be trusted,
and here it could not be.

Taken together, these points describe a method that would produce a lot of
numbers that could not be properly checked. For a bachelor's thesis, a dense
automatic result that cannot be verified is worse than a smaller set of values
that can. This is the reason I chose to measure the thickness by hand at a
limited number of well-imaged locations instead. Measuring by hand, however, only
works if the mesh is correctly placed inside the volume, so the first step before
measuring was to align the two. This is described next.

\section{Aligning the mesh and the volume}\label{sec:align}

The mesh and the volume did not overlap when they were first loaded together in both MeVisLab and 3D Slicer - only Blender could directly open the .usda file which preserved both coordinate systems in the correct orientation;
in MeVisLab and 3D Slicer they sat in different coordinate systems. This alignment matters here because, as
described in Section~\ref{sec:measure}, the thickness is measured along lines
perpendicular to the mesh surface. If the mesh is in the wrong place, those
perpendicular lines start in the wrong place, and the measurement is meaningless.
So the mesh and the volume had to be brought into the same coordinate frame
first.

Aligning the two by their shapes directly is difficult, because the leaflet
surface is smooth and curved and offers no clear matching landmarks to line up by
eye. Instead, I used a reference that all three programs agree on: the bounding
box. MeVisLab, 3D~Slicer and Blender each draw an axis-aligned bounding box
around the volume, and this box is the same in all of them because it is derived
from the same volume. Since everything was perfectly alligned in Blender I added 
planes to the mesh that correspond to the bounding
box of the volume, so that the mesh carries the same box as a built-in reference. Matching the
mesh's bounding planes to the volume's bounding box then defines the alignment
completely, because the two boxes have to coincide.

From this box-to-box correspondence a transformation matrix was obtained, which
was applied to the mesh in 3D~Slicer. After applying it, the mesh sits inside the
volume in the correct position and orientation, so that a leaflet in the mesh
lines up with the corresponding bright band in the scan. Using the bounding box
as a shared reference is more reliable here than trying to register the curved
surfaces against each other, and it keeps the alignment reproducible, because the
box is defined by the data and not by a subjective judgement of where the
surfaces should meet.

\section{Measuring leaflet thickness in 3D~Slicer}\label{sec:measure}

With the mesh aligned, the thickness was measured in 3D~Slicer \cite{fedorov20123d}. The tool used is
the line (ruler) markup, which gives a direct distance reading in millimetres
between two points. The important part is not the tool but the rule I followed
for placing the two points, because a clear rule is what makes the measurements
repeatable.

I chose not to judge the leaflet edges purely by eye, and I also chose not to
threshold the whole volume, for the reasons given in
Section~\ref{sec:recover}. Instead, I used a local intensity criterion, applied
only at the spot being measured. Tissue was defined as an intensity of ~0.2
or higher. Each measurement line was placed perpendicular to the mesh surface at
the chosen location. Along that line, the measurement starts where the intensity
first rises to~0.2 and ends where it falls back on or below~0.2, and the thickness is
the length of that span. This is the same kind of intensity criterion that the
automatic approach would have used, but applying it locally, along a single line
at a spot I had judged to be clean, avoids the problem of applying it to the whole
noisy volume at once. I chose where to place each line; the intensity rule then
decided where the tissue began and ended.

Where the measurements were taken was itself a choice driven by the data quality.
Thickness is not the same everywhere on a leaflet, so I wanted to sample across
the surface, but two regions turned out to be hard to measure. The first is the
centre of the valve, around the opening, where the probe left a noisy hole in the
scan. In that region it is difficult to see a clean tissue band, so reliable
perpendicular measurements are hard to obtain. The second is the thin outer brim
of the leaflet, at the edge of the mesh, where measurements were also difficult
to take cleanly. In between these two problem regions the data is clearer, so
most of the measurements were taken there, where the tissue band is well defined
and the intensity rule can be applied with confidence.

For the two difficult regions I made a deliberate simplification rather than
forcing measurements the data does not support. For the noisy centre I took all
the measurements I could obtain there and reduced them to a single value, the
median, and treated the thickness across the whole centre region as constant at
that value. I did the same for the outer edge, taking the median of the edge
measurements and treating the edge as having one constant thickness. The median
was used rather than the mean because it is less affected by the occasional bad
reading, which is exactly what one expects in these noisy regions. Assuming a
homogeneous thickness in the centre and at the edge is a simplification, and it is
stated as such, but it is a reasonable one: the data in those regions is not good
enough to justify anything finer, and a single well-chosen value is more honest
than a set of unreliable point measurements.

\section{Spreading the measurements over the mesh}\label{sec:interp}

At this point there is, for each leaflet, a set of thickness values: several point
measurements in the well-imaged middle region, plus one constant value for the
noisy centre and one constant value for the outer edge. The mesh, however, has
several thousand vertices, and the step that builds the solid
(Section~\ref{sec:solid}) needs a thickness value at every single vertex. So the
measured values have to be spread over the whole surface.

The measured points can be recorded as follows:

\begin{table}[htbp]
  \centering
  \begin{tabular}{clcc}
    \toprule
    Point & Location & Leaflet & Thickness [mm] \\
    \midrule
    1 & Base               & \ph{Mesh1} & \ph{} \\
    2 & Middle             & \ph{Mesh1} & \ph{} \\
    3 & Free edge (median) & \ph{Mesh1} & \ph{} \\
    4 & Centre (median)    & \ph{Mesh1} & \ph{} \\
    \dots & \dots          & \dots      & \dots \\
    \bottomrule
  \end{tabular}
  \caption{Measured leaflet thickness at the sampled locations.}
  \label{tab:measurements}
\end{table}

The simplest way to fill the gap would be to give the whole leaflet one average
thickness. I rejected this because it would remove the variation between the base
and the edge that I went to the trouble of capturing. So some form of
interpolation is needed for the region that is well sampled, while the centre and
the edge keep their constant median values as described in
Section~\ref{sec:measure}.

For the interpolation I chose inverse distance weighting (IDW) \cite{shepard1968two}. The idea is easy
to state: each vertex is given a thickness that is a weighted average of the known
values, and values that are closer to the vertex count more than ones further
away. The weight of each value falls off with distance, following one over the
distance raised to a power. A vertex sitting right next to a known value gets
almost exactly that value, and a vertex halfway between two of them gets roughly
their average.

I chose IDW over more advanced interpolation methods such as kriging or radial
basis functions on purpose. With only a handful of known values, the more
advanced methods bring in extra assumptions and parameters that could not be
checked with so little data. IDW, in contrast, makes no assumption about an
underlying function, works with irregularly placed values, and is simple enough
to explain and defend in full. For the power parameter, a value of two was used. This gives greater weight to nearby measurements while still providing a smooth transition between the measured points.

The interpolation was carried out in Blender using a short Python script, as shown in Listing~\ref{lst:thickness_interpolation}. NumPy \cite{harris2020array} is used for the distance and weighted averaging calculations. For each vertex, the script calculates the Euclidean distance to all measured points, determines the corresponding inverse-distance weights, and computes the weighted average of the measured thickness values. For every vertex the script computes the distances to
the known values, forms the weights, and takes the weighted average. The result
is stored on the mesh as a per-vertex weight, called a vertex group in Blender.
Before storing, the values are scaled to the range zero to one by dividing by the
largest thickness, because the modifier used in the next step expects weights in
that range.

\begin{lstlisting}[language=Python, caption={Interpolation of the thickness distribution in Blender using inverse distance weighting.}, label={lst:thickness_interpolation}]
import bpy
import numpy as np


def apply_thickness(obj_name, known_pos, known_val, power=2):
    obj = bpy.data.objects[obj_name]
    mesh = obj.data

    all_verts = np.array([v.co for v in mesh.vertices])

    def idw(query, kpos, kval, power):
        diff = query[:, None, :] - kpos[None, :, :]
        dist = np.maximum(np.linalg.norm(diff, axis=2), 1e-9)

        w = 1.0 / dist**power

        return (w * kval).sum(1) / w.sum(1)

    thickness_mm = idw(
        all_verts,
        np.array(known_pos),
        np.array(known_val),
        power
    )

    max_d = np.array(known_val).max()

    # Remove existing group for re-runs
    if "thickness" in obj.vertex_groups:
        obj.vertex_groups.remove(
            obj.vertex_groups["thickness"]
        )

    vg = obj.vertex_groups.new(name="thickness")

    for i, d in enumerate(thickness_mm):
        vg.add(
            [i],
            float(d / max_d),
            'REPLACE'
        )

    print(
        f"{obj_name}: "
        f"{thickness_mm.min():.2f} - "
        f"{thickness_mm.max():.2f} mm "
        f"(max={max_d})"
    )

# ------------------------------------------------------------------
# Mesh 1
# ------------------------------------------------------------------
apply_thickness(
    "mesh1",
    known_pos=[
        [point1x,  point1y,  point1z],
        [point2x,  point2y,  point2z],
        [point3x,  point3y,  point3z],
        [point4x,  point4y,  point4z],
        [point5x,  point5y,  point5z],
        [point6x,  point6y,  point6z],
        [point7x,  point7y,  point7z],
        [point8x,  point8y,  point8z],
        [point9x,  point9y,  point9z],
        [point10x, point10y, point10z],
        [point11x, point11y, point11z],
    ],
    known_val=[
        thickness1,
        thickness2,
        thickness3,
        thickness4,
        thickness5,
        thickness6,
        thickness7,
        thickness8,
        thickness9,
        thickness10,
        thickness11,
    ]
)

# ------------------------------------------------------------------
# Mesh 2
# ------------------------------------------------------------------
apply_thickness(
    "mesh2",
    known_pos=[
        [point1x, point1y, point1z],
        [point2x, point2y, point2z],
        [point3x, point3y, point3z],
        [point4x, point4y, point4z],
        [point5x, point5y, point5z],
        [point6x, point6y, point6z],
    ],
    known_val=[
        thickness1,
        thickness2,
        thickness3,
        thickness4,
        thickness5,
        thickness6,
    ]
)
\end{lstlisting}

One limitation of this choice should be stated openly, because it also explains a
later decision. IDW here uses straight-line (euclidean) distance through space,
not distance measured along the curved surface. Across a fold, or across the gap
between the two leaflets, straight-line distance could mix values that really
should not be mixed, because two points can be close in space while being far
apart along the tissue. This is exactly why the two leaflets are kept as two
separate meshes and interpolated separately: a value measured on one leaflet
cannot leak across the gap onto the other. Within a single, fairly flat leaflet
the difference between straight-line distance and along-surface distance is small
enough to accept for the purpose of this work.

\section{Building the solid model}\label{sec:solid}

The last step has this problem to solve: there is now a thickness value at every
vertex, and a single surface that represents the underside of the leaflet, and
this has to be turned into a closed solid whose wall has the correct local
thickness.

For this step, Blender's Solidify modifier was used \cite{blenderSolidify}. This modifier takes a
surface and extrudes it by a chosen thickness to produce a solid, and,
importantly, it can scale that thickness per vertex using a vertex group. This is
what makes it fit the situation here: the interpolated thickness field from
Section~\ref{sec:interp} is stored as exactly such a vertex group, so it maps
directly onto how far each part of the surface is pushed out. In practice the
modifier's thickness is set to the largest thickness value, given in metres
because Blender works in metres, and the zero-to-one weights then scale every
other part of the surface below that maximum. The local thickness at a vertex is
the weight at that vertex multiplied by the modifier thickness.

The extrusion is done in one direction only, which is where the check from
Section~\ref{sec:input} finally pays off. Because the mesh is the underside of
the leaflet and not the midsurface, the whole thickness has to be added on one
side, growing outward from the surface, rather than split half on each side. This
is also the concrete reason why the ``two times the distance'' rule from the
distance-field idea in Section~\ref{sec:recover} would have been wrong for this
data: a factor of two is only correct when the surface is the centre, and here it
is a boundary. In the modifier this is set by using a one-sided offset so that the
material grows only outward. The direction of the extrusion follows the surface
normals, so the normals were checked beforehand and flipped where needed, to make
sure the wall grows away from the leaflet body and not into it.

Each leaflet is solidified on its own before the two are combined. The reason is
that the two leaflets have different thickness ranges, and a single modifier can
only hold one thickness value and one vertex group. So \ph{Mesh1} and \ph{Mesh2}
are each given their own modifier with their own maximum thickness. Only after the
thickness has been fixed into the geometry, by applying the modifier, are the two
meshes joined into one object. Joining them earlier would force a single shared
thickness on both, and would also let the interpolation bleed across the gap
between them, for the reason given in Section~\ref{sec:interp}.

The output of this step is a single solid mesh in which the leaflet walls have the
measured, locally varying thickness.

\section{Printing the complete model}\label{sec:print}
Once a sufficiently accurate approximation, or model, of the mitral valve had been successfully developed, the actual objective and core of this work was addressed: the fabrication of a realistic ultrasound phantom. As accuracy and realism were prioritized over ease of production and scalability, FDM 3D printing was selected as the manufacturing method \cite{medellin2019design}. Modern 3D printers provide sufficient detail resolution and, importantly for this project, the ability to work with flexible materials such as TPU. The fabrication of flexible structures using molds and a material such as silicone would also have been possible and could even have enabled scalable production. However, 3D printing was ultimately chosen because it allows for easy modifications to the model and eliminates the need to consider how the mold can be removed after casting.

Throughout the project, while the thickness measurements were being determined, test prints were conducted in parallel to identify the approach that would yield the best possible result. It was apparent that the heart valve would require support structures during printing, as the convex geometry would need to rest as flat as possible on the build plate, resulting in overhangs of well above 60° \cite{medellin2019design}. Printing the model vertically slightly reduced the cross-sectional area requiring support, but not to a significant extent. Furthermore, this orientation caused the structure to become unstable during printing, as the build plate moves back and forth along the Y-axis. The higher the print progressed, the more pronounced this movement became, causing the structure to wobble.

A major disadvantage of support material is the difficulty of completely removing all residues and achieving a smooth surface after printing. Test prints using PLA, for example, which is a relatively rigid plastic, demonstrated that removing the support structures alone was already a considerable challenge. Achieving a surface as smooth as the upper side of the print would then require substantial additional effort. Printing with flexible TPU did not significantly change this situation. The TPA95A used was specified with a Shore A hardness of 95; in general, lower Shore A values correspond to softer and more flexible TPU materials \cite{recreusFilamentFilaflex90}. It could be printed using settings that were almost identical to those used for PLA. For TPU82A \cite{recreusFilamentFilaflex82}, the printing speed was reduced to avoid compromising print quality. However, at the desired thickness of the heart valve, this material remained too inflexible. TPU60A \cite{recreusFilamentFilaflex60}, which provided an appropriate degree of flexibility and softness, introduced another problem: the amount of support material used previously was insufficient to adequately support the material between the individual support structures. The resulting print contained numerous holes and areas where the printed material had fused with the supports. The solution was to increase the density of the support structures by reducing the spacing between them. In this case, this was achieved by selecting smaller squares in the grid pattern.

This resulted in a print using the desired material that appeared intact at first glance. However, the issue of the difficult-to-remove support structures remained. To address this problem, additional prototypes were produced using a multi-nozzle printer capable of printing with multiple materials simultaneously. The idea was to print the supports using a water-soluble filament, in this case polyvinyl alcohol (PVA) \cite{moradi2023optimizing}, while continuing to print the mitral valve using TPU60A. The finished model could then have been placed in water, allowing the supports to dissolve completely without leaving any residue.

However, there were concerns that the TPU might not adhere sufficiently to the PVA supports during printing, potentially resulting in print defects in the model. Therefore, PLA supports were initially used instead. Ultimately, this proved to be the right choice for this particular application, as the PLA and TPU adhered strongly enough to ensure a successful print while still allowing the valve to be easily peeled away from the support material after completion. As hoped, this resulted in a perfectly smooth underside. Furthermore, the gap within the valve opening was preserved.

The only remaining step was to manually create a small incision in the opening using a scalpel. This was necessary because 3D printing small openings generally leaves behind thin strands of filament, a phenomenon known as \textit{stringing}, which can cause the opening to become partially sealed. The completed prototype, including the support structures, was handed over to Siemens Healthineers, where the valve is currently undergoing further testing. \todo{change if Siemens gets back with results}